\chapter{Rings, Ideals, and Localization}

The integers have addition and multiplication, tied together by the distributive
laws. A ring retains exactly this structure. The familiar examples---integers,
residue classes, polynomials, and matrices---will give the definitions their
meaning. Ideals permit us to force selected elements to be zero; localization
instead forces selected elements to be invertible.

\section{Rings, Subrings, and Homomorphisms}
\begin{definition}
A \emph{ring} is a set \(R\) with an abelian-group structure under \(+\),
an associative multiplication, and a multiplicative identity \(1\), such
that, for all \(a,b,c\in R\),
\[
\begin{array}{rclcrcl}
(a+b)+c&=&a+(b+c),&&a+b&=&b+a,\\
a+0&=&a,&&a+(-a)&=&0,\\
(ab)c&=&a(bc),&&1a&=&a1=a,\\
a(b+c)&=&ab+ac,&&(a+b)c&=&ac+bc.
\end{array}
\]
A ring is \emph{commutative} if \(ab=ba\) for all \(a,b\). A subring
contains the ambient identity \(1_R\) and is closed under subtraction and
multiplication.
\end{definition}

If \(1=0\), then \(R\) has only one element: for every \(a\in R\),
\[
a=1a=0a=0.
\]
Thus the zero ring is included in the definition.

\begin{unremark}[Ring conventions]
Terminology varies among authors.  Some allow rings without an identity,
and some call a subset a subring when it is closed under subtraction and
multiplication without requiring it to contain the ambient identity.  Here
rings have identities, ring homomorphisms preserve them, and subrings contain
the same identity as the ambient ring.  For example, \(2\mathbb Z\) is
closed under subtraction and multiplication, so it is a subring of
\(\mathbb Z\) under the looser convention, but it is not a subring in these
notes because \(1\notin2\mathbb Z\).
\end{unremark}

The rings \(\mathbb Z,\mathbb Q,\mathbb R\), and a field \(F\), are
commutative. So are \(\mathbb Z/n\mathbb Z\), \(F[x]\), and \(R\times S\),
where operations are coordinatewise. Matrix multiplication makes \(M_n(F)\)
a ring, usually not a commutative one. Thus commutativity must be an
additional hypothesis, not a consequence of the axioms.

\begin{proposition}[Basic identities]
For \(a,b\in R\), one has
\[0a=a0=0,\qquad (-a)b=a(-b)=-(ab),\qquad(-a)(-b)=ab.\]
\end{proposition}
\begin{proof}
Since \(0a=(0+0)a=0a+0a\), additive cancellation gives \(0a=0\);
the proof of \(a0=0\) is analogous. Now
\(ab+(-a)b=(a+(-a))b=0\), giving the next identities. Applying this twice
gives \((-a)(-b)=-\bigl(a(-b)\bigr)=ab\).
\end{proof}

\begin{definition}
A \emph{unit} is an element \(u\) with \(uv=vu=1\) for some \(v\).
A nonzero \(a\) is a \emph{zero divisor} if \(ab=0\) or \(ba=0\) for some
nonzero \(b\). A \emph{division ring} is a nonzero ring in which every
nonzero element is a unit; commutativity is not assumed.  A \emph{field} is
a commutative division ring.  An \emph{integral domain} (or simply a
\emph{domain}) is a commutative ring with \(1\ne0\) and no zero divisors.
\end{definition}

Every field is a domain: if \(ab=0\) and \(a\ne0\), multiply by \(a^{-1}\).
The converse fails, since \(\mathbb Z\) is a domain but \(2\) has no inverse.
In \(\mathbb Z/12\mathbb Z\), \(5\) is a unit, while \(2\cdot6=0\) shows
that \(2\) and \(6\) are zero divisors. The class of \(a\) modulo \(n\) is
a unit exactly when \(\gcd(a,n)=1\), by B\'ezout's identity.

\begin{example}[Hamilton's quaternions]
The set
\[
\mathbb H=\{a+bi+cj+dk:a,b,c,d\in\mathbb R\}
\]
is a division ring under distributive multiplication determined by
\[
i^2=j^2=k^2=ijk=-1.
\]
For example, \(ij=k\), whereas \(ji=-k\); thus \(\mathbb H\) is not
commutative.  Define
\[
\overline{a+bi+cj+dk}=a-bi-cj-dk,
\qquad N(a+bi+cj+dk)=a^2+b^2+c^2+d^2.
\]
Expanding with the defining relations gives
\(q\overline q=\overline q q=N(q)\) and
\(N(qr)=N(q)N(r)\).  Hence, if \(q\ne0\), then \(N(q)>0\) and
\[
q^{-1}=\frac{\overline q}{N(q)}.
\]
The quaternions are therefore the standard example of a noncommutative
division ring.
\end{example}

\begin{theorem}
Every finite integral domain is a field.
\end{theorem}
\begin{proof}
Let \(0\ne a\in R\).  The map \(R\to R\), \(x\mapsto ax\), is injective:
\(ax=ay\) implies \(a(x-y)=0\), so cancellation gives \(x=y\).  Since
\(R\) is finite, the map is surjective.  Thus \(ax=1\) for some \(x\in R\),
and commutativity gives \(xa=1\).  Every nonzero element is invertible.
\end{proof}
\begin{applicationtheorem}[Wedderburn's little theorem]
Every finite division ring is a field.
\end{applicationtheorem}
\begin{remark}
This assertion is considerably deeper than the finite-domain theorem: the
latter uses cancellation and a finite-set argument, whereas the former must
first force commutativity.  Its proof is beyond the needs of the present
discussion.
\end{remark}

\begin{example}[A useful warning]
Over a domain, every unit of \(R[x]\) is constant. This is false in a ring
with nilpotents: in \((\mathbb Z/4\mathbb Z)[x]\),
\[(1+2x)^2=1+4x+4x^2=1.\]
Thus \(1+2x\) is a nonconstant unit. Domain hypotheses in polynomial
arguments are therefore essential.
\end{example}

\begin{definition}
A ring homomorphism \(f:R\to S\) preserves addition, multiplication, and
identity. Its kernel is \(\ker f=\{a:f(a)=0\}\), and its image is
\(\operatorname{im}f=\{f(a):a\in R\}\).
\end{definition}
For example, reduction \(\mathbb Z\to\mathbb Z/n\mathbb Z\) is onto and has
kernel \(n\mathbb Z\). As for group homomorphisms, \(f(0)=0\) and
\(f(-a)=-f(a)\). The determinant is not a ring homomorphism on matrices,
because it does not in general preserve addition.

\section{Ideals and Quotient Rings}
An additive subgroup is not automatically suitable for quotienting: it must
remain stable when multiplied by arbitrary ring elements. This is precisely
the ring analogue of the normal-subgroup condition.

\begin{definition}
A \emph{left ideal} of \(R\) is an additive subgroup \(I\) such that
\(ra\in I\) for every \(r\in R\), \(a\in I\).  A \emph{right ideal} is
defined by the condition \(ar\in I\).  A \emph{two-sided ideal} is both a
left and a right ideal.  Throughout these notes, unless explicitly
qualified, an \emph{ideal} means a two-sided ideal, and we write
\(I\trianglelefteq R\).

For \(A\subseteq R\), \((A)\) is the smallest two-sided ideal containing
\(A\).  If \(A=\{a_1,\ldots,a_m\}\), then
\[(a_1,\ldots,a_m)=
\left\{\sum_{\ell=1}^N r_\ell a_{i_\ell}s_\ell:
N\ge0,\ 1\le i_\ell\le m,\ r_\ell,s_\ell\in R\right\}.\]
When \(R\) is commutative, left, right, and two-sided ideals coincide, and
this simplifies to the familiar sums
\[
\sum_k r_ka_k.
\]
\end{definition}

The ideals \((0)\) and \((1)=R\) are trivial. In \(\mathbb Z\),
\((n)=n\mathbb Z\), while in \(F[x]\), \((x)\) is the set of polynomials
with zero constant term. In \(F[x,y]\), \((x^2,xy)\) consists of sums
\(Ax^2+Bxy\). The odd integers are not an ideal of \(\mathbb Z\), since
\(2\cdot1\) is not odd.  In \(M_n(F)\), the matrices whose last row is zero
form a right ideal but generally not a left ideal.  Quotient rings therefore
require two-sided, not merely one-sided, ideals.

The notation \((A)\) is deliberately parallel to the subgroup generated by
a subset: it records the smallest ideal forced by the elements of \(A\).
In a commutative ring, the principal ideal generated by \(a\) is simply
\((a)=\{ra:r\in R\}\).  Principal ideals reverse divisibility:
\[(a)\subseteq(b)\quad\Longleftrightarrow\quad b\mid a.\]
Indeed, \(a\in(b)\) is exactly the assertion \(a=br\).  Thus, in
\(\mathbb Z\), \((12)\subseteq(6)\subseteq(2)\).  This reversal will
underlie the use of principal ideals, gcds, ascending chains, and
factorization in Chapter~4.

\begin{proposition}
Every ideal of \(\mathbb Z\) is \(n\mathbb Z\) for a unique \(n\ge0\).
\end{proposition}
\begin{proof}
For a nonzero ideal \(I\), choose its least positive member \(n\). For
\(a\in I\), write \(a=qn+r\), \(0\le r<n\). Since \(r=a-qn\in I\),
minimality gives \(r=0\); hence \(I=n\mathbb Z\). The reverse containment
is ideal closure, and \(n\ge0\) makes the generator unique.
\end{proof}

For ideals \(I,J\), define their sum and product by
\[
I+J=\{a+b:a\in I,\ b\in J\},\qquad
IJ=\left\{\sum_{k=1}^n a_kb_k:a_k\in I,\ b_k\in J\right\}.
\]
The sums are finite; the empty sum is \(0\).  Together with \(I\cap J\),
these are ideals.  The sum is the smallest ideal containing both: any ideal
containing \(I\) and \(J\) contains every \(a+b\).  The intersection is their
largest common subideal. Always \(IJ\subseteq I\cap J\).  Finite sums and
products \(I_1+\cdots+I_n\) and \(I_1\cdots I_n\) are defined similarly.
For example, in \(\mathbb Z\),
\[(6)+(10)=(2),\qquad (6)\cap(10)=(30),\qquad (6)(10)=(60).\]
Thus intersection and product need not agree when the ideals are not
comaximal.

\begin{theorem}[Quotient Ring Criterion]
For an additive subgroup \(I\le R\), the rule
\[(a+I)(b+I)=ab+I\]
is well-defined exactly when \(I\) is a two-sided ideal. In that case \(R/I\)
is a ring and \(\pi(a)=a+I\) is a surjective homomorphism.
\end{theorem}
\begin{proof}
If \(a'=a+i\), \(b'=b+j\) with \(i,j\in I\), then
\[(a+i)(b+j)-ab=aj+ib+ij\in I,\]
using closure under multiplication on both sides.  Thus the product is
independent of representatives.  All ring axioms descend from \(R\).
Conversely, \(i+I=0+I\) implies
\((r+I)(i+I)=(r+I)(0+I)\), hence \(ri\in I\), while
\((i+I)(r+I)=(0+I)(r+I)\) gives \(ir\in I\). Thus \(I\) is two-sided.
\end{proof}

\begin{theorem}[Universal Property of Quotients]
If \(f:R\to T\) is a homomorphism and \(I\subseteq\ker f\), then there is a
unique homomorphism \(\bar f:R/I\to T\) such that \(f=\bar f\circ\pi\).
\end{theorem}
\[
\begin{tikzcd}[column sep=large, row sep=large]
R \arrow[r, "\pi"] \arrow[dr, "f"'] & R/I \arrow[d, "\bar f"] \\
& T
\end{tikzcd}
\]
\begin{proof}
Put \(\bar f(a+I)=f(a)\). If \(a+I=b+I\), then \(a-b\in I\), so
\(f(a)=f(b)\); this proves well-definedness. The homomorphism laws follow
from those for \(f\), and every coset is \(\pi(a)\), proving uniqueness.
\end{proof}

Thus \(R/I\) is universal among maps that force elements of \(I\) to zero,
just as quotient groups were in Part~I.  The structural parallel is
\[
\text{normal subgroup }N\trianglelefteq G\ \longleftrightarrow\ G/N,
\qquad
\text{two-sided ideal }I\trianglelefteq R\ \longleftrightarrow\ R/I.
\]
In particular, taking \(I=R\) gives \(R/R\), the zero ring.

\begin{theorem}[Isomorphism and correspondence]
For \(f:R\to S\), \(\ker f\) is an ideal and
\[R/\ker f\cong\operatorname{im}f.\]
Ideals of \(R/I\) correspond to ideals of \(R\) containing \(I\), by
\(J\mapsto\pi^{-1}(J)\) and \(K\mapsto K/I\); in particular,
\((R/I)/(K/I)\cong R/K\) if \(I\subseteq K\).
\end{theorem}
\begin{proof}
The kernel is additive and \(f(ra)=f(r)f(a)=0=f(ar)\) for \(a\in\ker f\).
The map \(a+\ker f\mapsto f(a)\) is well-defined, onto the image, and has
zero kernel. Directly checking the two displayed assignments shows they
are inverse ideal-preserving maps. Finally,
\((r+I)+(K/I)\mapsto r+K\) is well-defined and has zero kernel.
\end{proof}

\section{Characteristic and the Prime Subring}
Let \(R\) be a ring with identity. The canonical homomorphism
\[
\iota:\mathbb Z\longrightarrow R,\qquad n\longmapsto n1_R
\]
records how the integers act inside \(R\). Its kernel is an ideal of
\(\mathbb Z\), so it is either \(\{0\}\) or \(n\mathbb Z\) for a unique
integer \(n>0\).

\begin{definition}\label{def:ring-characteristic}
If \(\ker\iota=\{0\}\), the \emph{characteristic} of \(R\) is \(0\). If
\(\ker\iota=n\mathbb Z\) with \(n>0\), its characteristic is \(n\). The image
of \(\iota\) is the \emph{prime subring} of \(R\).
\end{definition}

Under this definition the zero ring has characteristic \(1\), since
\(1_R=0\) and therefore \(\ker\iota=\mathbb Z\).

The First Isomorphism Theorem gives
\[
\operatorname{im}\iota\cong\mathbb Z/\ker\iota.
\]
Thus in characteristic \(0\) the prime subring is isomorphic to
\(\mathbb Z\), while in characteristic \(p>0\) it is
\(\mathbb Z/p\mathbb Z\), and hence isomorphic to \(\mathbb F_p\) when
\(p\) is prime.

\begin{proposition}[Characteristic of a domain]
\label{prop:domain-characteristic}
If \(R\) is a nonzero integral domain, then
\(\operatorname{char}R\) is either \(0\) or a prime.
\end{proposition}
\begin{proof}
Suppose \(\operatorname{char}R=n>0\). If \(n=ab\) with
\(1<a,b<n\), then minimality of \(n\) gives
\[
a1_R\ne0,\qquad b1_R\ne0,
\]
whereas
\[
(a1_R)(b1_R)=n1_R=0.
\]
This contradicts the absence of zero divisors. Hence \(n\) is prime.
\end{proof}

For example,
\[
\operatorname{char}\mathbb Z=\operatorname{char}\mathbb Q=0,\qquad
\operatorname{char}\mathbb F_p=p,\qquad
\operatorname{char}(\mathbb Z/n\mathbb Z)=n\quad(n\geq2).
\]

\begin{example}[Group rings]
Let \(R\) be a commutative ring and \(G\) a group.  The \emph{group ring}
\(R[G]\) consists of finite formal sums
\[
\sum_{g\in G}a_g g,
\]
with coefficientwise addition and multiplication determined by distributivity
and the multiplication in \(G\):
\[
\left(\sum_g a_g g\right)\left(\sum_h b_h h\right)
=\sum_{g,h}a_gb_h(gh).
\]
The finiteness condition makes each coefficient in the product a finite sum.
For the cyclic group \(C_2=\{1,s\}\), \(s^2=1\), every element is uniquely
\(a+bs\), and the map sending \(s\) to the class of \(x\) gives
\[
R[C_2]\cong R[x]/(x^2-1).
\]
This is an elementary way in which a group relation becomes a polynomial
relation.  The polynomial-ring construction is developed formally in the next
chapter; this is an intentional forward preview.  No representation theory is
needed here.
\end{example}

\section{Prime Ideals, Maximal Ideals, Domains, and Fields}
From this point through localization, rings are assumed to be commutative
and unital; the zero ring remains allowed.  Thus the unqualified word
\emph{ideal} continues to mean a two-sided ideal, but sidedness is automatic.
\begin{definition}
A proper ideal \(\mathfrak p\) is \emph{prime} if
\(ab\in\mathfrak p\) implies \(a\in\mathfrak p\) or \(b\in\mathfrak p\).
A proper ideal \(\mathfrak m\) is \emph{maximal} if no ideal is strictly
between \(\mathfrak m\) and \(R\).
\end{definition}
\begin{theorem}\label{thm:prime-maximal-quotient}
\(\mathfrak p\) is prime iff \(R/\mathfrak p\) is a domain, and
\(\mathfrak m\) is maximal iff \(R/\mathfrak m\) is a field. Hence every
maximal ideal is prime.
\end{theorem}
\begin{proof}
\((a+\mathfrak p)(b+\mathfrak p)=0\) iff \(ab\in\mathfrak p\), giving the
first claim. If \(\mathfrak m\) is maximal and \(a\notin\mathfrak m\), then
\((a)+\mathfrak m=R\), so \(ra+m=1\) for suitable \(r,m\), and
\(r+\mathfrak m\) inverts \(a+\mathfrak m\). Conversely, a field has only
the ideals \(0\) and itself; apply the correspondence theorem.
\end{proof}
In \(\mathbb Z\), \((p)\) is prime and maximal for prime \(p\), while
\((6)\) is neither: \(2\cdot3\in(6)\), but \(2,3\notin(6)\). In
\(\mathbb Z/12\mathbb Z\), the ideal \((2)\) is
maximal because its quotient is \(\mathbb Z/2\mathbb Z\). Later, the
division algorithm will show that \((f)\subset F[x]\) is maximal exactly
when \(f\) is irreducible.

\begin{applicationtheorem}[Maximal ideals above proper ideals]
\label{thm:maximal-ideal-existence}
Let \(R\) be a commutative ring with identity. Every proper ideal \(I\) of
\(R\) is contained in a maximal ideal.
\end{applicationtheorem}
\begin{proof}
Let \(\mathcal P\) be the set of proper ideals of \(R\) that contain \(I\),
ordered by inclusion. It is nonempty because \(I\in\mathcal P\). If
\(\mathcal C\) is a chain in \(\mathcal P\), then
\[
J=\bigcup_{K\in\mathcal C}K
\]
is an ideal containing \(I\): the chain condition puts any two elements of
\(J\) in a common member of \(\mathcal C\), where the ideal operations may
be performed. Moreover \(J\) is proper. Indeed, if \(1\in J\), then
\(1\in K\) for some \(K\in\mathcal C\), contradicting the propriety of
\(K\). Thus every chain has an upper bound in \(\mathcal P\). Zorn's Lemma
now gives a maximal member \(\mathfrak m\) of \(\mathcal P\), which is
maximal among all proper ideals of \(R\), hence is a maximal ideal.
\end{proof}
\begin{unremark}[A set-theoretic qualification]
The preceding proof invokes Zorn's Lemma. For Noetherian rings, the
corresponding maximality conclusion follows from the usual ascending-chain
maximality argument, so no appeal to Zorn's Lemma is needed in the standard
treatment. We do not pursue the precise foundational strength of the general
maximal-ideal theorem here.
\end{unremark}

\section{The Chinese Remainder Theorem}
\begin{definition}
Ideals \(I,J\) of a commutative ring are \emph{comaximal} if \(I+J=R\).
Equivalently,
\[
I+J=R\quad\Longleftrightarrow\quad 1=a+b
\quad\text{for some }a\in I,\ b\in J.
\]
Indeed, \(1\in I+J\) is exactly the displayed assertion, and an ideal
containing \(1\) is all of \(R\).
\end{definition}
\begin{lemma}
If \(I+J=R\), then \(I\cap J=IJ\).
\end{lemma}
\begin{proof}
Certainly \(IJ\subseteq I\cap J\). Choose \(u\in I,v\in J\) with
\(u+v=1\). If \(x\in I\cap J\), then \(x=xu+xv\), and both summands lie in
\(IJ\).
\end{proof}
\begin{theorem}[Chinese Remainder Theorem]
If \(I+J=R\), then \(R/(I\cap J)\cong R/I\times R/J\).
\end{theorem}
\begin{proof}
The homomorphism \(\Phi(r)=(r+I,r+J)\) has kernel \(I\cap J\). With
\(u+v=1\), \(u\in I,v\in J\), the element \(av+bu\) maps to
\((a+I,b+J)\); hence \(\Phi\) is onto. Apply the First Isomorphism Theorem.
\end{proof}
\begin{example}
To solve \(x\equiv2\pmod3\), \(x\equiv3\pmod5\), use
\(6+(-5)=1\), with \(6\in(3)\), \(-5\in(5)\). The proof gives
\(x=2(-5)+3(6)=8\). All solutions are \(8\pmod {15}\).
\end{example}
\begin{theorem}[Finite Chinese Remainder Theorem]
If \(I_1,\ldots,I_n\) are pairwise comaximal ideals of a commutative ring
\(R\), then
\[
R/\bigcap_{i=1}^n I_i\cong\prod_{i=1}^n R/I_i.
\]
Moreover,
\[
\bigcap_{i=1}^n I_i=\prod_{i=1}^n I_i.
\]
\end{theorem}
\begin{proof}
For \(n=2\), both assertions were just proved.  Suppose they hold for
\(I_1,\ldots,I_{n-1}\), and put \(J=I_1\cap\cdots\cap I_{n-1}\).  For each
\(i<n\), choose \(a_i\in I_i\), \(b_i\in I_n\) with \(a_i+b_i=1\).
Then \(\prod_{i<n}a_i\in J\), while
\(1-\prod_{i<n}a_i\in I_n\), by expanding the product.  Thus \(J+I_n=R\).
The two-ideal theorem and the induction hypothesis now give
\[
R/(J\cap I_n)\cong R/J\times R/I_n
\cong\prod_{i=1}^n R/I_i.
\]
The two-ideal product formula gives
\(J\cap I_n=JI_n\), and induction identifies this with
\(I_1\cdots I_n\).
\end{proof}

\section{Localization and Local Rings}
The inclusion \(\mathbb Z\subseteq\mathbb Q\) permits division by every
nonzero integer.  Often we want only selected denominators.  For instance,
\[\mathbb Z[1/2]=\{a/2^n:a\in\mathbb Z,\ n\ge0\}\]
allows division by powers of \(2\), without requiring arbitrary rational
denominators.  Localization is the controlled construction that introduces
chosen denominators: it makes selected elements units while otherwise
changing the ring as little as the required universal property permits.
The canonical map need not be injective when zero divisors are present.

If \(s\in R\), the basic case is
\[R[1/s]=\{1,s,s^2,\ldots\}^{-1}R.\]
Its elements are represented by \(a/s^n\), and \(s\) becomes a unit with
inverse \(1/s\).  Equivalently, \(R[1/s]\) is universal among rings receiving
a homomorphism from \(R\) in which the image of \(s\) is invertible.
\(\mathbb Z[1/2]\) is the first example.  Inverting \(6\) also makes both
\(2\) and \(3\) invertible, so one should not expect the only new units to
be powers of the chosen element.

\begin{definition}
A set \(S\subseteq R\) is \emph{multiplicatively closed} if \(1\in S\) and
\(s,t\in S\) imply \(st\in S\). Define \((a,s)\sim(b,t)\) in \(R\times S\)
when \(u(at-bs)=0\) for some \(u\in S\). The class of \((a,s)\) is \(a/s\),
and \(S^{-1}R\) has operations
\[\frac as+\frac bt=\frac{at+bs}{st},\qquad
\frac as\frac bt=\frac{ab}{st}.\]
\end{definition}
Ordinary fraction arithmetic suggests
\[
\frac as=\frac bt\quad\Longleftrightarrow\quad at=bs.
\]
In a general ring, however, a denominator may be a zero divisor, so the
cancellation implicit in cross multiplication can fail after denominators
are inverted.  The condition \(u(at-bs)=0\) for some \(u\in S\) says exactly
that the cross-difference becomes zero after multiplying by an allowed
denominator; it is therefore the correct replacement.  If every element of
\(S\) is a non-zero-divisor, this condition simplifies to \(at=bs\), since
one may cancel \(u\).  The conditions in the definition are also forced by
fraction arithmetic: if \(s,t\) are allowed denominators, then
\((a/s)(b/t)=ab/(st)\) requires \(st\) to be allowed as well; \(1\) is needed
to represent each original element as \(a/1\).
\begin{proposition}
The relation is an equivalence relation, and these operations are
well-defined.
\end{proposition}
\begin{proof}
Reflexivity uses \(1(as-as)=0\), and symmetry changes the sign. If
\(u(at-bs)=0\), \(v(bq-ct)=0\), then multiplying the first equality by
\(vq\), the second by \(us\), and adding gives \(uv(atq-cst)=0\).
This proves transitivity. If representatives are changed, multiplying the
two witnessing equalities by the other numerator and denominator proves
\(ab/(st)=a'b'/(s't')\). For sums the needed cross-difference is
\(tt'(as'-a's)+ss'(bt'-b't)\), after multiplication by the two witnesses;
it is zero. The remaining ring laws follow from common denominators.
\end{proof}
\begin{proposition}[Fraction fields as localizations]
If \(R\) is an integral domain and \(S=R\setminus\{0\}\), then
\[\operatorname{Frac}(R)=S^{-1}R.\]
Under this identification, \(a/b=c/d\) if and only if \(ad=bc\).
In particular, \(\mathbb Q=\operatorname{Frac}(\mathbb Z)\).
\end{proposition}
\begin{proof}
The set \(S\) is multiplicatively closed, and every one of its elements is
regular.  The defining equivalence relation is consequently ordinary cross
multiplication.  Thus the usual fractions and the localization have the same
representatives, equality relation, and operations.  Taking \(R=\mathbb Z\)
gives the final assertion.
\end{proof}
\begin{example}
In \(R=\mathbb Z/6\mathbb Z\), take \(S=\{1,3\}\). Then
\(0/1=2/1\), since \(3(0-2)=0\), although \(0\ne2\). The witness \(u\) is
therefore essential in rings with zero divisors.  It also shows directly that
the canonical map \(R\to S^{-1}R\) need not be injective.
\end{example}
\begin{proposition}
The canonical map \(\iota:R\to S^{-1}R\), \(a\mapsto a/1\), has kernel
\[\{a\in R:sa=0\text{ for some }s\in S\}.\]
It is injective exactly when every \(s\in S\) is \emph{regular}, meaning
\(sa=0\Rightarrow a=0\).  In the present commutative setting, this says
precisely that \(S\) contains no zero divisors.  Moreover, \(s/1\) is a unit
with inverse \(1/s\).
\end{proposition}
\begin{proof}
\(a/1=0/1\) is exactly the displayed condition. The fraction formulas
prove the homomorphism statement and \((s/1)(1/s)=1\).
\end{proof}
\begin{theorem}[Universal Property of Localization]
If \(f:R\to T\) maps each \(s\in S\) to a unit, there is a unique
\(\widetilde f:S^{-1}R\to T\) such that \(f=\widetilde f\circ\iota\).
\end{theorem}
\[
\begin{tikzcd}[column sep=large, row sep=large]
R \arrow[r, "\iota"] \arrow[dr, "f"'] & S^{-1}R \arrow[d, "\widetilde f"] \\
& T
\end{tikzcd}
\]
\begin{proof}
Set \(\widetilde f(a/s)=f(a)f(s)^{-1}\). A witnessing equality
\(u(at-bs)=0\), after applying \(f\) and cancelling the unit \(f(u)\),
proves this is well-defined. The fraction formulas prove the homomorphism
laws. Finally \(a/s=(a/1)(s/1)^{-1}\), which forces uniqueness.
\end{proof}
Quotients force elements of \(I\) to be zero; localizations force elements
of \(S\) to be units. These are complementary universal constructions.  In
particular, \(R/(s)\) forces \(s\) to be zero, whereas \(R[1/s]\) forces
\(s\) to be invertible.  For example,
\(\mathbb Z/(2)=\mathbb F_2\), while \(\mathbb Z[1/2]\) still has
characteristic zero.

\begin{definition}
A commutative ring is \emph{local} if it has a unique maximal ideal.
\end{definition}
\begin{proposition}
A commutative ring with identity is local if and only if its nonunits form an
ideal.
\end{proposition}
\begin{proof}
If \(\mathfrak m\) is the unique maximal ideal, every element outside it is a
unit: otherwise its principal ideal is proper and, by the
\hyperref[thm:maximal-ideal-existence]{maximal-ideal theorem}, lies in a
maximal ideal, which must be \(\mathfrak m\). Hence the nonunits are exactly
\(\mathfrak m\). Conversely, suppose that the nonunits form an ideal \(N\).
It is proper because \(1\) is a unit. Every proper ideal is contained in
\(N\), since an ideal containing a unit is all of \(R\). Thus \(N\) itself
is maximal, and every maximal ideal, being proper, equals \(N\). This proves
both existence and uniqueness.
\end{proof}
\begin{unremark}
Under this convention the zero ring is not local: it has no maximal ideals.
\end{unremark}

For a prime ideal \(\mathfrak p\), \(1\notin\mathfrak p\), and its complement
\(S=R\setminus\mathfrak p\) is multiplicatively closed.  Indeed, if
\(a,b\notin\mathfrak p\) but \(ab\in\mathfrak p\), primality would put
\(a\) or \(b\) in \(\mathfrak p\), a contradiction.  Hence
\[R_{\mathfrak p}=(R\setminus\mathfrak p)^{-1}R\]
is defined, and its elements are fractions \(a/s\) with \(s\notin\mathfrak p\).
It makes every element outside \(\mathfrak p\) invertible, leaving the
behavior attached to \(\mathfrak p\) visible; informally, it focuses attention
at \(\mathfrak p\).
\begin{theorem}
For a prime ideal \(\mathfrak p\), \(R_{\mathfrak p}\) is local with unique
maximal ideal
\[\mathfrak pR_{\mathfrak p}=\{a/s:a\in\mathfrak p,\ s\notin\mathfrak p\}.\]
\end{theorem}
\begin{proof}
\(R\setminus\mathfrak p\) is multiplicatively closed by primality. The
displayed set is an ideal: numerators in \(\mathfrak p\) remain there under
addition and multiplication. It is proper, since an equality \(1=a/s\)
with \(a\in\mathfrak p\) would say \(u(s-a)=0\) for some
\(u\notin\mathfrak p\), contradicting primality. If \(a\notin\mathfrak p\),
then \(a/s\) has inverse \(s/a\). If \(a\in\mathfrak p\), it cannot be a
unit by the same argument applied to \((a/s)(b/t)=1\). Hence this ideal is
exactly the nonunits; every proper ideal is contained in it.
\end{proof}
For \(\mathbb Z_{(p)}\), the unique maximal ideal is
\(p\mathbb Z_{(p)}\).  Concretely,
\[\mathbb Z_{(p)}=\{a/b\in\mathbb Q:p\nmid b\};\]
every integer not divisible by \(p\) is a unit, and its unique maximal ideal
is formed by the fractions whose numerator is divisible by \(p\).

\[
\begin{array}{c|c|c}
S&\text{localization}&\text{elements made invertible}\\ \hline
\{1,s,s^2,\ldots\}&R[1/s]&\text{one chosen element }s\\
R\setminus\mathfrak p&R_{\mathfrak p}&\text{everything outside }\mathfrak p\\
R\setminus\{0\}\ (R\text{ a domain})&\operatorname{Frac}(R)&
\text{every nonzero element}
\end{array}
\]

\section*{Exercises}
\begin{exercise}[Basic] Determine the units and zero divisors of
\(\mathbb Z/10\mathbb Z\). \end{exercise}
\begin{exercise}[Core] Prove that the preimage of an ideal under a ring
homomorphism is an ideal. Use correspondence to describe ideals of
\(\mathbb Z/18\mathbb Z\). \end{exercise}
\begin{exercise}[Core] Show \((x)\) and \((x-1)\) are comaximal in \(F[x]\),
and describe \(F[x]/(x(x-1))\). \end{exercise}
\begin{exercise}[Core] Determine which of \((2),(3),(4),(6)\) are prime or
maximal in \(\mathbb Z/12\mathbb Z\). \end{exercise}
\begin{exercise}[Further] Construct \(S^{-1}\mathbb Z\cong\mathbb Z[1/2]\)
for \(S=\{1,2,2^2,\ldots\}\), and identify its units. \end{exercise}
\begin{exercise}[Further] Prove
\(\mathbb Z_{(p)}/p\mathbb Z_{(p)}\cong\mathbb F_p\). \end{exercise}
